\documentclass[10pt,a4paper]{article}
\usepackage[margin=20mm,headheight=15pt,headsep=7mm,footskip=10mm]{geometry}
\usepackage{fontspec,amsmath,amssymb,mathtools,microtype}
\setmainfont{lmroman10-regular.otf}[BoldFont=lmroman10-bold.otf,ItalicFont=lmroman10-italic.otf,BoldItalicFont=lmroman10-bolditalic.otf]
\setsansfont{lmsans10-regular.otf}[BoldFont=lmsans10-bold.otf,ItalicFont=lmsans10-oblique.otf,BoldItalicFont=lmsans10-boldoblique.otf]
\setmonofont[Scale=.86]{lmmono10-regular.otf}[ItalicFont=lmmono10-italic.otf]
\usepackage{ragged2e}\usepackage{xcolor,booktabs,longtable,array,graphicx,enumitem,fancyhdr,needspace,titlesec,hyperref,xurl}
\definecolor{AtlasNavy}{HTML}{18334C}\definecolor{AtlasTeal}{HTML}{116D73}\definecolor{AtlasGray}{HTML}{52616D}
\hypersetup{colorlinks=true,linkcolor=AtlasTeal,urlcolor=AtlasTeal,pdfcreator={Complexity Atlas},unicode=true}
\urlstyle{same}\setlength{\parindent}{0pt}\setlength{\parskip}{4.8pt}\setlength{\emergencystretch}{3em}
\renewcommand{\baselinestretch}{1.08}\setcounter{secnumdepth}{0}\setlist{nosep,leftmargin=1.5em}
\titleformat{\section}{\Large\bfseries\sffamily\color{AtlasNavy}\raggedright}{}{0em}{}
\titleformat{\subsection}{\large\bfseries\sffamily\color{AtlasTeal}\raggedright}{}{0em}{}
\titleformat{\subsubsection}{\normalsize\bfseries\color{AtlasNavy}\raggedright}{}{0em}{}
\titlespacing*{\section}{0pt}{14pt}{6pt}\titlespacing*{\subsection}{0pt}{10pt}{5pt}
\pagestyle{fancy}\fancyhf{}\fancyhead[L]{\small\sffamily\color{AtlasGray}COMPLEXITY ATLAS}
\fancyhead[R]{\small\sffamily\color{AtlasGray}ENGLISH EDITION}\fancyfoot[C]{\small\thepage}\renewcommand{\headrulewidth}{.3pt}
\newcommand{\AtlasMath}[1]{\par\begingroup\centering\setbox0=\hbox{$\displaystyle #1$}\ifdim\wd0>1.4\linewidth\typeout{ATLAS WIDE FORMULA: \the\wd0}\fi\ifdim\wd0>\linewidth\resizebox{\linewidth}{!}{\box0}\else\box0\fi\par\endgroup}
\newcommand{\AtlasInlineMath}[1]{\begingroup\setbox0=\hbox{$#1$}\ifdim\wd0>.4\linewidth\endgroup\AtlasMath{#1}\else\endgroup$#1$\fi}
\pdfstringdefDisableCommands{\def\AtlasInlineMath#1{#1}}
\DeclareMathOperator{\polylog}{polylog}\DeclareMathOperator{\dist}{dist}\DeclareMathOperator{\diam}{diam}
\DeclareMathOperator{\rank}{rank}\DeclareMathOperator{\tr}{tr}\DeclareMathOperator{\diag}{diag}
\DeclareMathOperator{\prox}{prox}\DeclareMathOperator{\KL}{KL}\DeclareMathOperator{\Var}{Var}
\providecommand{\eps}{\varepsilon}\providecommand{\R}{\mathbb R}\providecommand{\E}{\mathbb E}
\providecommand{\mathbbm}[1]{\mathbf{#1}}\providecommand{\mathds}[1]{\mathbf{#1}}
\providecommand{\llbracket}{[\![}\providecommand{\rrbracket}{]\!]}
\begin{document}
\raggedright

{\LARGE\bfseries\sffamily\color{AtlasNavy} Nonconvex stationarity with separate exact gradient and Hessian queries under joint smoothness\par}\vspace{5pt}


\textbf{Setting ID.} \texttt{SO\_\allowbreak{}HO\_\allowbreak{}HOI09\_\allowbreak{}S\_\allowbreak{}joint\_\allowbreak{}L1L2}



Single-objective optimization



The strict gradient-Hessian budget frontier is determined for derivative-only information under simultaneous gradient and Hessian Lipschitz bounds.



\textbf{Status.} Fully optimal on the documented comparison domain



\textbf{TASK tags.} solution concept: Gradient stationarity; task family: Single objective



\textbf{MODEL tags.} smoothness: Lipschitz gradient; smoothness: Lipschitz Hessian



\textbf{ORACLE tags.} feedback quality: Exact; operation: Gradient; operation: Hessian; public information: Hidden objects explicitly declared; public information: Public inputs explicitly declared; scalar derivative order: Scalar-function first derivative; scalar derivative order: Scalar-function second derivative



\textbf{ALGORITHMS tags.} randomization: Deterministic; restriction: No additional span restriction explicitly imposed



\textbf{OUTPUT tags.} postprocessing: Unqueried output explicitly allowed



\textbf{GUARANTEE tags.} aggregation: Per instance; quantifier: Algorithm before instance



\textbf{RESOURCE tags.} budget: Fixed worst-case budget; measure: N gradient; measure: N hessian



\needspace{7\baselineskip}\section{1. Setting specification}


\needspace{4\baselineskip}\subsection{Task}


Find \AtlasInlineMath{\widehat x} satisfying \AtlasInlineMath{\|\nabla f(\widehat x)\|_2\le\varepsilon} while minimizing the gradient-query cap under a prescribed strict full-Hessian-query cap.



\needspace{4\baselineskip}\subsection{Admissible model}


For each public finite dimension \AtlasInlineMath{d}, the unknown \AtlasInlineMath{f\in C^2(\mathbb R^d)} has finite infimum, \AtlasInlineMath{f(x_0)-\inf f\le\Delta}, \AtlasInlineMath{\operatorname{Lip}(\nabla f)\le L_1} and \AtlasInlineMath{\operatorname{Lip}(\nabla^2f)\le L_2}. All budgets and the initial point are public; derivative norms are induced Euclidean operator norms. There is no convexity or query-radius restriction.



\needspace{4\baselineskip}\subsection{Oracle interface}


At an arbitrary point choose exactly one of \AtlasInlineMath{O_g(x)=\nabla f(x)} and \AtlasInlineMath{O_H(x)=\nabla^2f(x)}. A Hessian query returns the full matrix but neither the function value nor its gradient. The two counters are charged separately, and there is no function-value channel.



\needspace{4\baselineskip}\subsection{Output}


The output may be any transcript-dependent point of \AtlasInlineMath{\mathbb R^d}, including an unqueried point.



\needspace{4\baselineskip}\subsection{Loss or accuracy criterion}


The error is \AtlasInlineMath{\|\nabla f(\widehat x)\|_2} and the guarantee must hold deterministically on every instance.



\needspace{4\baselineskip}\subsection{Resource accounting}


A feasible pair \AtlasInlineMath{(g,h)\in\mathbb N_0^2} is witnessed by one algorithm with \AtlasInlineMath{N_g\le g} and \AtlasInlineMath{N_H\le h} on every input. Define \AtlasInlineMath{G^\star_\infty(h;\varepsilon)} by minimizing \AtlasInlineMath{g} with a guarantee uniform over all finite dimensions. Internal subproblem solves and exact-real arithmetic are uncharged.



\needspace{4\baselineskip}\subsection{Algorithm views and success guarantees}


Queries, stopping decisions and outputs are arbitrary measurable functions of the observed history and public inputs. No linear-span or zero-respecting restriction is imposed. Internal exact-real computation and storage are uncharged.



The algorithm is deterministic and uses no private random seed.




\AtlasMath{\forall f:\ \|\nabla f(\widehat x)\|_2\le\varepsilon}




Success and both resource caps hold on every instance.



\textbf{Budget convention.} fixed\_worst\_case\_cap



\needspace{7\baselineskip}\section{2. Notation and parameters}


\AtlasInlineMath{d}: Public finite ambient dimension \textbf{Domain:} \AtlasInlineMath{d\in\mathbb N_{\ge1}}.



\AtlasInlineMath{L_1}: Global gradient Lipschitz bound \textbf{Domain:} \AtlasInlineMath{L_1>0}.



\AtlasInlineMath{L_2}: Global Hessian Lipschitz bound \textbf{Domain:} \AtlasInlineMath{L_2>0}.



\AtlasInlineMath{\Delta}: Initial objective gap \textbf{Domain:} \AtlasInlineMath{\Delta>0}.



\AtlasInlineMath{\varepsilon}: Gradient-norm tolerance \textbf{Domain:} \AtlasInlineMath{\varepsilon>0}.



\AtlasInlineMath{h}: Strict integer Hessian-query cap \textbf{Domain:} \AtlasInlineMath{h\in\mathbb N_0}.



\AtlasInlineMath{A}: Second-order complexity scale \textbf{Domain:} \AtlasInlineMath{A>0}.



\AtlasInlineMath{A=\Delta\sqrt{L_2}\varepsilon^{-3/2}}



\AtlasInlineMath{B}: Gradient-only complexity scale \textbf{Domain:} \AtlasInlineMath{B>0}.



\AtlasInlineMath{B=\Delta\sqrt{L_1}L_2^{1/4}\varepsilon^{-7/4}}



\needspace{12\baselineskip}\section{3. Complexity bracket}


\needspace{4\baselineskip}\subsection{Compared quantities}


\needspace{4\baselineskip}\subsubsection{Q1 - minimax metric resource quantity}


\AtlasMath{G^\star_\infty(h;\varepsilon)}


\textbf{Scope.} The union of the explicit parameter and dimension domains of the linked comparison regimes; each bound retains its own applicability conditions.



\textbf{Optimization direction.} minimize



\textbf{Resource contract.} expression: \AtlasInlineMath{(N_g,N_H),\ N_H\le h;\ \min N_g}; description: The Hessian cap is an integer hard cap; a rich derivative does not include a gradient.



\needspace{4\baselineskip}\subsection{R1 - No Hessian calls}


\textbf{Region.} \AtlasInlineMath{0<\varepsilon\le c_0\min\{L_1^2/L_2,\Delta^{2/3}L_2^{1/3}\}} for the sufficiently small universal \AtlasInlineMath{c_0} in equation regime; hence \AtlasInlineMath{B\ge A\ge128}. \AtlasInlineMath{h=0}. Upper bounds hold in every finite dimension. The Hessian-rich product lower bound is one-dimensional. For the scarce-Hessian lower bound, Lemma matrix-lift gives sufficient dimension \AtlasInlineMath{C(Jd_0+J+g+h+1)}, with \AtlasInlineMath{J=2(h+1)}, \AtlasInlineMath{d_0\le C_{\rm dim}m} from Lemma local-gadget, and local scales \AtlasInlineMath{\delta=\Delta/(6J)}, \AtlasInlineMath{\nu=L_2/C_*}, \AtlasInlineMath{\ell=\min\{L_1/C_*,\delta\nu/\varepsilon\}} from equation poor-calibration-final. The construction dimension is allowed to depend on the proposed pair; no fixed-dimensional matching claim is made.



\textbf{Quantity.} Q1



\textbf{Resource.} The Hessian cap is an integer hard cap; a rich derivative does not include a gradient.



\textbf{Status.} Optimal



\textbf{Lower bound.} B2



\AtlasMath{\Omega\!\left(B\right)}


\textbf{Upper bound.} B1



\AtlasMath{O\!\left(B\right)}


\textbf{Gap or comparison.} Constant-factor matching throughout this branch.



\needspace{4\baselineskip}\subsection{R2 - Scarce Hessian budget}


\textbf{Region.} \AtlasInlineMath{0<\varepsilon\le c_0\min\{L_1^2/L_2,\Delta^{2/3}L_2^{1/3}\}} for the sufficiently small universal \AtlasInlineMath{c_0} in equation regime; hence \AtlasInlineMath{B\ge A\ge128}. \AtlasInlineMath{h\in\mathbb N}, \AtlasInlineMath{1\le h\le A}. Upper bounds hold in every finite dimension. The Hessian-rich product lower bound is one-dimensional. For the scarce-Hessian lower bound, Lemma matrix-lift gives sufficient dimension \AtlasInlineMath{C(Jd_0+J+g+h+1)}, with \AtlasInlineMath{J=2(h+1)}, \AtlasInlineMath{d_0\le C_{\rm dim}m} from Lemma local-gadget, and local scales \AtlasInlineMath{\delta=\Delta/(6J)}, \AtlasInlineMath{\nu=L_2/C_*}, \AtlasInlineMath{\ell=\min\{L_1/C_*,\delta\nu/\varepsilon\}} from equation poor-calibration-final. The construction dimension is allowed to depend on the proposed pair; no fixed-dimensional matching claim is made.



\textbf{Quantity.} Q1



\textbf{Resource.} The Hessian cap is an integer hard cap; a rich derivative does not include a gradient.



\textbf{Status.} Optimal



\textbf{Lower bound.} B4



\AtlasMath{\Omega\!\left(\min\{B,A^{3/2}/\sqrt h\}\right)}


\textbf{Upper bound.} B3



\AtlasMath{O\!\left(\min\{B,A^{3/2}/\sqrt h\}\right)}


\textbf{Gap or comparison.} Constant-factor matching throughout this branch.



\needspace{4\baselineskip}\subsection{R3 - Hessian-rich budget}


\textbf{Region.} \AtlasInlineMath{0<\varepsilon\le c_0\min\{L_1^2/L_2,\Delta^{2/3}L_2^{1/3}\}} for the sufficiently small universal \AtlasInlineMath{c_0} in equation regime; hence \AtlasInlineMath{B\ge A\ge128}. \AtlasInlineMath{h\in\mathbb N}, \AtlasInlineMath{h>A}. Upper bounds hold in every finite dimension. The Hessian-rich product lower bound is one-dimensional. For the scarce-Hessian lower bound, Lemma matrix-lift gives sufficient dimension \AtlasInlineMath{C(Jd_0+J+g+h+1)}, with \AtlasInlineMath{J=2(h+1)}, \AtlasInlineMath{d_0\le C_{\rm dim}m} from Lemma local-gadget, and local scales \AtlasInlineMath{\delta=\Delta/(6J)}, \AtlasInlineMath{\nu=L_2/C_*}, \AtlasInlineMath{\ell=\min\{L_1/C_*,\delta\nu/\varepsilon\}} from equation poor-calibration-final. The construction dimension is allowed to depend on the proposed pair; no fixed-dimensional matching claim is made.



\textbf{Quantity.} Q1



\textbf{Resource.} The Hessian cap is an integer hard cap; a rich derivative does not include a gradient.



\textbf{Status.} Optimal



\textbf{Lower bound.} B6



\AtlasMath{\Omega\!\left(\max\{1,A^2/h\}\right)}


\textbf{Upper bound.} B5



\AtlasMath{O\!\left(\max\{1,A^2/h\}\right)}


\textbf{Gap or comparison.} Constant-factor matching throughout this branch.



\needspace{7\baselineskip}\section{4. Bounds and supporting evidence}


The following conditions apply to every bound in this profile, in addition to each bound's stated region and algorithm scope.



\textbf{Dimension conditions.} Upper bounds hold in every finite dimension. The Hessian-rich product lower bound is one-dimensional. For the scarce-Hessian lower bound, Lemma matrix-lift gives sufficient dimension \AtlasInlineMath{C(Jd_0+J+g+h+1)}, with \AtlasInlineMath{J=2(h+1)}, \AtlasInlineMath{d_0\le C_{\rm dim}m} from Lemma local-gadget, and local scales \AtlasInlineMath{\delta=\Delta/(6J)}, \AtlasInlineMath{\nu=L_2/C_*}, \AtlasInlineMath{\ell=\min\{L_1/C_*,\delta\nu/\varepsilon\}} from equation poor-calibration-final. The construction dimension is allowed to depend on the proposed pair; no fixed-dimensional matching claim is made.



\needspace{4\baselineskip}\subsection{B1 - Gradient-only method and the imported first-order lower bound}


\textbf{Admission.} Included in the bracket



\textbf{Direction.} upper



\textbf{Witness or lower-bound scope.} Deterministic witness



\AtlasMath{O\!\left(B\right)}


\textbf{Sources.} \href{https://github.com/WU-Wendao/Colombo/blob/main/colombo_sensorium/Complexity_altas/data/sources/colombo/0904/manuscript.pdf}{[1]}



\textbf{Conditions.} \AtlasInlineMath{0<\varepsilon\le c_0\min\{L_1^2/L_2,\Delta^{2/3}L_2^{1/3}\}} for the sufficiently small universal \AtlasInlineMath{c_0} in equation regime; hence \AtlasInlineMath{B\ge A\ge128}. \AtlasInlineMath{h=0}.; One algorithm attains both stated coordinate caps.



\textbf{Constants.} All unmarked constants are universal.



\textbf{Logarithmic factors.} None.



\textbf{Source location.} \href{https://github.com/WU-Wendao/Colombo/blob/main/colombo_sensorium/Complexity_altas/data/sources/colombo/0904/manuscript.pdf}{[1]} Corollary complete-frontier; Lemma gradient-only; verified labels: cor:complete-frontier (line 407), eq:complete-frontier (line 409), lem:gradient-only (line 644)



Gradient-only method and the imported first-order lower bound. Deterministic adaptive algorithms.



One algorithm attains both stated coordinate caps.



\needspace{4\baselineskip}\subsection{B2 - Gradient-only method and the imported first-order lower bound}


\textbf{Admission.} Included in the bracket



\textbf{Direction.} lower



\textbf{Witness or lower-bound scope.} Deterministic witness



\AtlasMath{\Omega\!\left(B\right)}


\textbf{Sources.} \href{https://github.com/WU-Wendao/Colombo/blob/main/colombo_sensorium/Complexity_altas/data/sources/colombo/0904/manuscript.pdf}{[1]}



\textbf{Conditions.} \AtlasInlineMath{0<\varepsilon\le c_0\min\{L_1^2/L_2,\Delta^{2/3}L_2^{1/3}\}} for the sufficiently small universal \AtlasInlineMath{c_0} in equation regime; hence \AtlasInlineMath{B\ge A\ge128}. \AtlasInlineMath{h=0}.; Arbitrary query points and unqueried outputs are allowed.



\textbf{Constants.} All unmarked constants are universal.



\textbf{Logarithmic factors.} None.



\textbf{Source location.} \href{https://github.com/WU-Wendao/Colombo/blob/main/colombo_sensorium/Complexity_altas/data/sources/colombo/0904/manuscript.pdf}{[1]} Corollary complete-frontier and its proof; verified labels: cor:complete-frontier (line 407), eq:complete-frontier (line 409)



Gradient-only method and the imported first-order lower bound. All deterministic measurable adaptive algorithms.



Arbitrary query points and unqueried outputs are allowed.



\needspace{4\baselineskip}\subsection{B3 - Strict integer lazy-extragradient schedule and sequential full-matrix resisting construction}


\textbf{Admission.} Included in the bracket



\textbf{Direction.} upper



\textbf{Witness or lower-bound scope.} Deterministic witness



\AtlasMath{O\!\left(\min\{B,A^{3/2}/\sqrt h\}\right)}


\textbf{Sources.} \href{https://github.com/WU-Wendao/Colombo/blob/main/colombo_sensorium/Complexity_altas/data/sources/colombo/0904/manuscript.pdf}{[1]}



\textbf{Conditions.} \AtlasInlineMath{0<\varepsilon\le c_0\min\{L_1^2/L_2,\Delta^{2/3}L_2^{1/3}\}} for the sufficiently small universal \AtlasInlineMath{c_0} in equation regime; hence \AtlasInlineMath{B\ge A\ge128}. \AtlasInlineMath{h\in\mathbb N}, \AtlasInlineMath{1\le h\le A}.; One algorithm attains both stated coordinate caps.



\textbf{Constants.} All unmarked constants are universal.



\textbf{Logarithmic factors.} None.



\textbf{Source location.} \href{https://github.com/WU-Wendao/Colombo/blob/main/colombo_sensorium/Complexity_altas/data/sources/colombo/0904/manuscript.pdf}{[1]} Theorem main part 2; Theorem strict; verified labels: thm:main (line 375), thm:strict (line 615)



Strict integer lazy-extragradient schedule and sequential full-matrix resisting construction. Deterministic adaptive algorithms.



One algorithm attains both stated coordinate caps.



\needspace{4\baselineskip}\subsection{B4 - Strict integer lazy-extragradient schedule and sequential full-matrix resisting construction}


\textbf{Admission.} Included in the bracket



\textbf{Direction.} lower



\textbf{Witness or lower-bound scope.} Deterministic witness



\AtlasMath{\Omega\!\left(\min\{B,A^{3/2}/\sqrt h\}\right)}


\textbf{Sources.} \href{https://github.com/WU-Wendao/Colombo/blob/main/colombo_sensorium/Complexity_altas/data/sources/colombo/0904/manuscript.pdf}{[1]}



\textbf{Conditions.} \AtlasInlineMath{0<\varepsilon\le c_0\min\{L_1^2/L_2,\Delta^{2/3}L_2^{1/3}\}} for the sufficiently small universal \AtlasInlineMath{c_0} in equation regime; hence \AtlasInlineMath{B\ge A\ge128}. \AtlasInlineMath{h\in\mathbb N}, \AtlasInlineMath{1\le h\le A}.; Arbitrary query points and unqueried outputs are allowed.



\textbf{Constants.} All unmarked constants are universal.



\textbf{Logarithmic factors.} None.



\textbf{Source location.} \href{https://github.com/WU-Wendao/Colombo/blob/main/colombo_sensorium/Complexity_altas/data/sources/colombo/0904/manuscript.pdf}{[1]} Theorem poor-complete; Corollary complete-frontier; verified labels: cor:complete-frontier (line 407), eq:complete-frontier (line 409), thm:poor-complete (line 1162)



Strict integer lazy-extragradient schedule and sequential full-matrix resisting construction. All deterministic measurable adaptive algorithms.



Arbitrary query points and unqueried outputs are allowed.



\needspace{4\baselineskip}\subsection{B5 - Gradient transport through Hessian integration and one-dimensional product obstruction}


\textbf{Admission.} Included in the bracket



\textbf{Direction.} upper



\textbf{Witness or lower-bound scope.} Deterministic witness



\AtlasMath{O\!\left(\max\{1,A^2/h\}\right)}


\textbf{Sources.} \href{https://github.com/WU-Wendao/Colombo/blob/main/colombo_sensorium/Complexity_altas/data/sources/colombo/0904/manuscript.pdf}{[1]}



\textbf{Conditions.} \AtlasInlineMath{0<\varepsilon\le c_0\min\{L_1^2/L_2,\Delta^{2/3}L_2^{1/3}\}} for the sufficiently small universal \AtlasInlineMath{c_0} in equation regime; hence \AtlasInlineMath{B\ge A\ge128}. \AtlasInlineMath{h\in\mathbb N}, \AtlasInlineMath{h>A}.; One algorithm attains both stated coordinate caps.



\textbf{Constants.} All unmarked constants are universal.



\textbf{Logarithmic factors.} None.



\textbf{Source location.} \href{https://github.com/WU-Wendao/Colombo/blob/main/colombo_sensorium/Complexity_altas/data/sources/colombo/0904/manuscript.pdf}{[1]} Theorem main part 3; Theorem transport; verified labels: thm:main (line 375), sec:transport (line 438), thm:transport (line 494)



Gradient transport through Hessian integration and one-dimensional product obstruction. Deterministic adaptive algorithms.



One algorithm attains both stated coordinate caps.



\needspace{4\baselineskip}\subsection{B6 - Gradient transport through Hessian integration and one-dimensional product obstruction}


\textbf{Admission.} Included in the bracket



\textbf{Direction.} lower



\textbf{Witness or lower-bound scope.} Deterministic witness



\AtlasMath{\Omega\!\left(\max\{1,A^2/h\}\right)}


\textbf{Sources.} \href{https://github.com/WU-Wendao/Colombo/blob/main/colombo_sensorium/Complexity_altas/data/sources/colombo/0904/manuscript.pdf}{[1]}



\textbf{Conditions.} \AtlasInlineMath{0<\varepsilon\le c_0\min\{L_1^2/L_2,\Delta^{2/3}L_2^{1/3}\}} for the sufficiently small universal \AtlasInlineMath{c_0} in equation regime; hence \AtlasInlineMath{B\ge A\ge128}. \AtlasInlineMath{h\in\mathbb N}, \AtlasInlineMath{h>A}.; Arbitrary query points and unqueried outputs are allowed.



\textbf{Constants.} All unmarked constants are universal.



\textbf{Logarithmic factors.} None.



\textbf{Source location.} \href{https://github.com/WU-Wendao/Colombo/blob/main/colombo_sensorium/Complexity_altas/data/sources/colombo/0904/manuscript.pdf}{[1]} Theorem lower; Corollary complete-frontier; verified labels: cor:complete-frontier (line 407), eq:complete-frontier (line 409), thm:lower (line 761)



Gradient transport through Hessian integration and one-dimensional product obstruction. All deterministic measurable adaptive algorithms.



Arbitrary query points and unqueried outputs are allowed.



\needspace{7\baselineskip}\section{Maintenance appendix: Native proof extensions}


Read-only presentation of the stored native record. Integrity hashes identify source values; this appendix grants no new theorem, source-read credit, review or approval.



Native pointer: /proofs; canonical SHA-256: 4f53cda18c2baa0c0354bb5f\allowbreak{}9a3ecbe5ed12ab4d8e11ba87\allowbreak{}3c2f11161202b945



Native value (empty\_array): []



\needspace{7\baselineskip}\section{References}


[1] Colombo\_0904. Sharp Deterministic Gradient--Hessian Tradeoffs for Finding Stationary Points. 0904\_gradient\_hessian\_tradeoffs.tex.



\url{https://github.com/WU-Wendao/Colombo/blob/main/colombo_sensorium/Complexity_altas/data/sources/colombo/0904/manuscript.pdf}

\end{document}
